\part{Computation}
\label{part:computation}

In Part 1: Fundamentals we have defined the basic quantities that are addressed in LQFT. Here we turn to numerical implementations of LQFTs and discuss the ways in which physical results are inferred from numerical calculations.
 First however, let us pause for a philosophical aside on the nature of computational quantum field theory. 
 
 Numerical computations by their very nature, and the nature of their implementations, deal with finite amounts of information. Renormalisable QFT on the other-hand deals with infinities, fields are defined at arbitrary small separations and often on non-compact manifolds, so mechanically specifying a quantum field by writing down its value at points in spacetime requires infinite information. Even in 0+1 dimensional field theories, Hilbert space is  infinite dimensional leading to the same conclusion. Progress is made by truncating this infinite system to a finite one and approaching the infinite system through a sequence of successively less-constraining truncations. As with classical numerical approaches, LQFT proceeds by discretising and compactifying spacetime through introducing the spacetime lattice and applying boundary conditions. State representation is also necessarily truncated by using finite-precision arithmetic or by more intentional Hilbert space truncation methods (see \ref{ch:quantum-simulation}). [LINK TO EFFECTIVE FIELD THEORY??] A numerical campaign may then evaluate quantities at different discretisation scales and in different space-time volumes (and in double vs single precision for example) and use these to understand the effects of the truncations. In the case of discretisation, the Symanzik improvement program and perturbation theory around fixed points of RG flows allows for prediction of the functional form of an extrapolation to the continuum limit. Similarly, effective field theories such as chiral perturbation theory constrain the approach to infinite volume in many contexts. However, we must \textit{assume} that the numerical results are in the regime of applicability of these theoretically-motivated extrapolations and test for their consistency. Numerical calculations show that LQCD with  Wilson action at $\beta\in[3,5]$ cannot be used to guide a continuum extrapolation at $\beta\to\infty$ due to a bulk transition. However, numerical calculations  and asymptotic scaling can not rule out the presence of another such transition at $\beta=305$ that invalidates the extensive calculations that are performed for $\beta\in[6,8]$ and used to make many predictions. While this is an assumption, significant evidence supports its validity, in particular while the Wilson gauge action has a bulk transition at $\beta\sim5.2$, improved actions and QCD with different types of fermions have very different bulk phase structures, yet show convergence in the continuum limit. The apparent universality of the continuum theory and  RG-based  arguments  about it make this particular scenario implausible, but do not rule it out.

 Very much like with QFT itself, whose nature is mathematically poorly defined but able to make exquisite predictions about the physical world, the purist may pause here and decide computational QFT is not for them. This is an entirely reasonable point of view and may lead you to new and fundamentally more powerful insights into QFT, but for those who are willing to accept certain assumptions in their scientific explorations, the numerical tools discussed below open vast realms of physics that we currently have no other way to address in such a systematically improvable manner. 

 Two classes of computational QFT are discussed in the following chapters. First is the Euclidean Monte-Carlo approach whose success has led to the dominance of LQCD in many aspects of nonperturbative QCD phenomenology (see Part \ref{part:pheno}). Algorithms and the techniques needed to extract physics from the results they produce are both introduced along with a  discussion of the contemporary role of machine learning methods in LQFT. Subsequently, we devote discussion to more exploratory approaches involving tensor networks and quantum computation as potential avenues to address quantities that are inaccessible to direct Monte-Carlo approaches, for example real-time dynamics and dense matter.
